\documentclass[../DoAn.tex]{subfiles}
\begin{document}
% Chương này có độ dài từ 9 đến 11 trang.

% Với phương pháp phân tích thiết kế hướng đối tượng, sinh viên sử dụng biểu đồ use case theo hướng dẫn của template này. Với các phương pháp khác, sinh viên trao đổi với giáo viên hướng dẫn để đổi tên và sắp xếp lại đề mục cho phù hợp. Ví dụ, thay vì sử dụng biểu đồ use case, sinh viên đi theo hướng tiếp cận Agile có thể dùng User Story.


% \textbf{Lưu ý}: Mỗi chương nên có thêm 1 đoạn mở đầu chương và kết thúc chương, mở đầu giới thiệu những nội dung sẽ trình bày trong chương, kết thúc tổng kết lại các nội dung đã trình bày

Chương 2 trình bày kết quả khảo sát hiện trạng và phân tích yêu cầu của hệ thống EduSchedule. Trước tiên, chương khảo sát phương pháp lập thời khóa biểu bằng bảng tính và một số phần mềm hỗ trợ hiện có. Tiếp theo, các quy trình nghiệp vụ chính của hệ thống được mô tả từ khâu chuẩn bị thông tin, phân công giảng dạy đến lập, công bố và tra cứu thời khóa biểu. Trên cơ sở đó, chương xác định các tác nhân, xây dựng biểu đồ use case, đặc tả những use case trọng tâm và trình bày các yêu cầu phi chức năng của hệ thống.
\section{Khảo sát hiện trạng}
\label{section:2.1}

Hiện nay, việc xây dựng thời khóa biểu tại các trường học có thể được thực hiện bằng nhiều phương pháp khác nhau, từ sử dụng bảng tính đến các phần mềm xếp thời khóa biểu chuyên dụng. Mỗi phương pháp có những ưu điểm và hạn chế riêng về khả năng quản lý thông tin, kiểm tra xung đột, tự động hóa quá trình sắp xếp và mức độ phù hợp với quy trình của từng nhà trường. Để xác định các yêu cầu cần thiết cho đề tài, tác giả tiến hành tìm hiểu thực trạng lập thời khóa biểu tại trường tiểu học và khảo sát một số giải pháp đang được sử dụng hiện nay.

Một trong những phương pháp thường được sử dụng để xây dựng thời khóa biểu là Microsoft Excel. Với phương pháp này, người phụ trách dựa trên danh sách lớp học, giáo viên, môn học, số tiết yêu cầu và kết quả phân công giảng dạy để sắp xếp từng tiết vào bảng thời khóa biểu. Sau khi hoàn thành, thời khóa biểu được rà soát, điều chỉnh và phổ biến tới giáo viên cùng các lớp học. Microsoft Excel có ưu điểm là quen thuộc, linh hoạt và cho phép người phụ trách chủ động điều chỉnh từng vị trí trong thời khóa biểu. Tuy nhiên, bảng tính không trực tiếp quản lý mối quan hệ giữa lớp học, môn học, giáo viên và phòng chức năng. Việc kiểm tra lịch dạy của giáo viên, số tiết đã được sắp xếp của từng môn và khả năng sử dụng phòng chức năng chủ yếu được thực hiện thủ công. Vì vậy, người phụ trách phải đối chiếu nhiều thông tin trong quá trình lập lịch và có thể bỏ sót xung đột, đặc biệt khi một giáo viên bộ môn được phân công giảng dạy tại nhiều lớp. Khi có thay đổi về giáo viên, phân công giảng dạy hoặc kế hoạch học tập, những vị trí liên quan trong thời khóa biểu phải được rà soát và điều chỉnh lại. Một thay đổi ở một lớp có thể ảnh hưởng đến lịch dạy của giáo viên tại các lớp khác. Bên cạnh đó, việc sử dụng nhiều bảng hoặc tệp riêng để quản lý thông tin có thể gây khó khăn trong quá trình cập nhật và kiểm tra tính thống nhất của dữ liệu.

Bên cạnh phương pháp sử dụng bảng tính, hiện nay đã có một số hệ thống hỗ trợ quản lý và xếp thời khóa biểu như VietSchool TKB, K12Online, TKB.com.vn và FTKB. Các sản phẩm này cung cấp nhiều chức năng như quản lý thông tin giảng dạy, hỗ trợ xếp lịch, kiểm tra xung đột, tra cứu thời khóa biểu và xuất dữ liệu. So với phương pháp thủ công, việc sử dụng phần mềm chuyên dụng giúp giảm một phần khối lượng kiểm tra và hỗ trợ người dùng quản lý thời khóa biểu tập trung hơn.

Các hệ thống như VietSchool và K12Online được phát triển trong phạm vi nền tảng quản lý giáo dục với nhiều nhóm chức năng khác nhau. Việc tích hợp thời khóa biểu với các nghiệp vụ quản lý nhà trường tạo thuận lợi cho những đơn vị đang sử dụng đồng bộ hệ sinh thái của sản phẩm. Tuy nhiên, phạm vi chức năng rộng cũng có thể khiến quy trình sử dụng trở nên phức tạp đối với trường chỉ có nhu cầu quản lý và lập thời khóa biểu ở quy mô nhỏ.

Trong khi đó, TKB.com.vn và FTKB tập trung nhiều hơn vào chức năng lập thời khóa biểu và hỗ trợ tự động hóa quá trình sắp xếp. Các hệ thống này có khả năng xử lý nhiều điều kiện của bài toán lập lịch và giúp người dùng tạo thời khóa biểu nhanh hơn so với phương pháp thủ công. Tuy nhiên, mỗi sản phẩm được xây dựng theo mô hình dữ liệu và quy trình riêng, vì vậy khả năng tùy chỉnh theo nghiệp vụ cụ thể của từng trường phụ thuộc vào những chức năng mà sản phẩm cung cấp.

Qua quá trình khảo sát, có thể thấy phương pháp sử dụng Excel đem lại sự chủ động trong việc sắp xếp nhưng thiếu khả năng quản lý dữ liệu tập trung và kiểm tra ràng buộc tự động. Ngược lại, các phần mềm chuyên dụng hỗ trợ nhiều chức năng tự động hóa nhưng có thể chưa phù hợp hoàn toàn với nhu cầu kết hợp giữa sắp xếp thủ công và tự động của người phụ trách tại trường tiểu học.

Từ những nhận xét trên, đề tài xác định nhu cầu xây dựng một hệ thống web tập trung vào quy trình lập thời khóa biểu cho trường tiểu học. Hệ thống cần hỗ trợ quản lý các thông tin liên quan, phân công giảng dạy, kiểm tra xung đột và hiển thị thời khóa biểu dưới dạng lưới trực quan. Người phụ trách có thể chủ động sắp xếp trước một số tiết học, sau đó sử dụng chức năng tự động để phân bổ các tiết còn thiếu mà không làm thay đổi các tiết đã có. Kết quả xếp tự động cần được hiển thị để người phụ trách kiểm tra và xác nhận trước khi lưu. Ngoài ra, hệ thống cần hỗ trợ giáo viên tra cứu thời khóa biểu được công bố và cho phép xuất thời khóa biểu ra file Excel.
\section{Mô tả nghiệp vụ}
\label{section:2.2}
Hệ thống được xây dựng nhằm hỗ trợ nhà trường trong quá trình chuẩn bị dữ liệu, phân công giảng dạy, xây dựng, điều chỉnh và công bố thời khóa biểu cho các lớp tiểu học. Phạm vi tin học hóa tập trung vào hoạt động xếp thời khóa biểu và khâu thiết lập các dữ liệu đầu vào phục vụ quá trình này, bao gồm dữ liệu về năm học, môn học, giáo viên, lớp học, phòng học chức năng và phân công giảng dạy. Trên cơ sở các dữ liệu đã được thiết lập, hệ thống hỗ trợ người phụ trách xây dựng, điều chỉnh và hoàn thiện thời khóa biểu trước khi đưa vào sử dụng.

Về danh mục môn học, khi người dùng đăng ký tài khoản, hệ thống tự động khởi tạo danh sách môn học theo từng khối lớp từ khối 1 đến khối 5, dựa trên Chương trình giáo dục phổ thông cấp tiểu học ban hành kèm theo Thông tư số 32/2018/TT-BGDĐT \cite{thongtu32} bao gồm các thông tin sau: tên môn học và định mức số tiết mỗi tuần theo từng khối. Đây là dữ liệu mặc định được sử dụng chung cho các năm học và học kỳ trong phạm vi tài khoản; người dùng có thể bổ sung, chỉnh sửa hoặc loại bỏ môn học, đồng thời điều chỉnh số tiết để phù hợp với kế hoạch giáo dục thực tế của nhà trường.

Về thông tin giáo viên, hệ thống quản lý tập trung danh sách giáo viên được sử dụng chung qua các năm học, bao gồm họ tên, loại giáo viên, các môn học có thể giảng dạy, trạng thái hoạt động và định mức số tiết tham khảo mỗi tuần. Người dùng có thể bổ sung hoặc cập nhật thông tin giáo viên; định mức số tiết được sử dụng để theo dõi và cảnh báo khối lượng giảng dạy trong quá trình phân công, không phải là ràng buộc bắt buộc khi xếp thời khóa biểu.

Về phòng chức năng, hệ thống quản lý danh sách các phòng phục vụ những môn học cần không gian riêng, bao gồm tên phòng, môn học liên kết và số lượng phòng hiện có. Người dùng có thể bổ sung, chỉnh sửa hoặc loại bỏ phòng chức năng để phù hợp với điều kiện cơ sở vật chất của nhà trường. Dữ liệu này được sử dụng khi xếp thời khóa biểu nhằm tránh bố trí số lớp sử dụng phòng tại cùng một thời điểm vượt quá số lượng phòng hiện có.

Về năm học, người phụ trách tạo năm học mới trên hệ thống. Mỗi năm học gồm hai học kỳ, trong đó học kỳ I có 18 tuần và học kỳ II có 17 tuần. Hệ thống quản lý một thời khóa biểu cho từng học kỳ; mỗi tuần lưu trữ lịch học của tất cả các lớp và có thể kế thừa thời khóa biểu từ tuần trước nếu không phát sinh thay đổi.

Trong từng năm học, người phụ trách thiết lập danh sách lớp theo từng khối từ khối 1 đến khối 5. Danh sách lớp được quản lý riêng theo năm học để phù hợp với kế hoạch tổ chức của nhà trường.

Sau khi đã có danh sách lớp học và danh sách giáo viên, người phụ trách thực hiện phân công giảng dạy, bao gồm phân công giáo viên chủ nhiệm và phân công giáo viên bộ môn. Đối với phân công chủ nhiệm, người phụ trách lựa chọn giáo viên chủ nhiệm cho từng lớp. Đối với phân công giáo viên bộ môn, người phụ trách lựa chọn giáo viên phụ trách từng môn học của từng lớp. Mỗi môn học tại một lớp chỉ do một giáo viên phụ trách; khi lựa chọn giáo viên mới, phân công trước đó được thay thế. Một giáo viên có thể được phân công giảng dạy nhiều môn hoặc nhiều lớp khác nhau tùy theo kế hoạch của nhà trường. 
Dựa trên định mức số tiết của các môn học và thông tin phân công giảng dạy, hệ thống tính tổng số tiết được phân công cho từng giáo viên. Khi tổng số tiết được phân công vượt quá định mức tham khảo, hệ thống hiển thị cảnh báo để người phụ trách xem xét và điều chỉnh khi cần thiết.

Thông tin lớp học, môn học, định mức số tiết của từng môn, giáo viên được phân công, phòng học và khung thời gian học là dữ liệu đầu vào trực tiếp cho quá trình xây dựng thời khóa biểu.

Hệ thống hỗ trợ xây dựng thời khóa biểu theo hai hình thức là sắp xếp thủ công và hỗ trợ sắp xếp tự động. Trong quá trình sắp xếp thủ công, người dùng lựa chọn môn học để bố trí vào từng ô tiết học tương ứng trên lưới thời khóa biểu, đồng thời có thể thay đổi hoặc loại bỏ môn học đã bố trí khi cần điều chỉnh. Khi sử dụng chức năng sắp xếp tự động, hệ thống giữ nguyên các tiết đã được bố trí thủ công và tự động sắp xếp các tiết học còn lại vào những vị trí phù hợp trên thời khóa biểu.

Trong quá trình chỉnh sửa thủ công, thời khóa biểu có thể tạm thời chưa đáp ứng đầy đủ các quy tắc nghiệp vụ, chẳng hạn còn thiếu tiết, xuất hiện khoảng trống trong buổi học hoặc phát sinh xung đột về giáo viên. Hệ thống phát hiện và thông báo các vi phạm để người dùng tiếp tục điều chỉnh. Việc cho phép tồn tại trạng thái tạm thời chưa hợp lệ giúp người dùng linh hoạt hơn trong quá trình xây dựng thời khóa biểu.

Trước khi thực hiện sắp xếp tự động, hệ thống kiểm tra tính đầy đủ của dữ liệu đầu vào, đặc biệt là việc phân công giáo viên cho các môn học của từng lớp. Nếu vẫn còn môn học chưa được phân công giáo viên, hệ thống không thực hiện sắp xếp tự động và yêu cầu người dùng hoàn thiện dữ liệu cần thiết.
Khi thực hiện sắp xếp tự động, hệ thống sử dụng thông tin về lớp học, môn học, số tiết cần bố trí, giáo viên được phân công, khung thời gian học và phòng học để tạo thời khóa biểu. Hệ thống bố trí các tiết học sao cho đáp ứng các quy tắc nghiệp vụ bắt buộc đã được xác định.

Một thời khóa biểu chỉ được coi là hợp lệ khi đáp ứng toàn bộ các ràng buộc bắt buộc, bao gồm:

\begin{itemize}[noitemsep]
    \item Mỗi lớp chỉ được học một môn tại cùng một tiết học.
    \item Một giáo viên không được giảng dạy nhiều lớp tại cùng một thời điểm.
    \item Mỗi môn học của từng lớp phải được bố trí đúng và đủ số tiết theo định mức của khối lớp.
    \item Các tiết học trong cùng một buổi của một lớp phải được bố trí liên tục, không xuất hiện tiết trống giữa các tiết học.
    \item Nếu một lớp được bố trí học vào buổi chiều thì các tiết học của buổi sáng trong cùng ngày phải được bố trí đầy đủ trước.
    \item Một phòng học, đặc biệt là phòng học chuyên biệt, không được sử dụng đồng thời cho nhiều lớp tại cùng một thời điểm.
    \item Mỗi tiết học phải được bố trí đúng giáo viên đã được phân công giảng dạy môn học của lớp.
    \item Các tiết học phải nằm trong khung thời gian và số tiết cho phép của từng buổi học.
\end{itemize}

Khi cần điều chỉnh trong quá trình giảng dạy, người dùng có thể lựa chọn áp dụng thay đổi từ một tuần xác định. Thời khóa biểu của các tuần trước thời điểm áp dụng vẫn được giữ nguyên, còn các tuần từ thời điểm được chọn trở đi được cập nhật theo thay đổi mới.

Trong quá trình chỉnh sửa, hệ thống cho phép lưu thời khóa biểu ở trạng thái tạm thời chưa hợp lệ để người dùng có thể tiếp tục hoàn thiện. Tuy nhiên, thời khóa biểu chỉ được xác nhận và công bố sau khi tất cả các vi phạm ràng buộc bắt buộc đã được xử lý.

Trong quá trình thực hiện kế hoạch giảng dạy, người dùng có thể thay đổi thông tin phân công giáo viên hoặc định mức số tiết của môn học. Khi dữ liệu đầu vào thay đổi, hệ thống kiểm tra lại thời khóa biểu có liên quan để xác định các vi phạm có thể phát sinh.

Đối với trường hợp thay đổi giáo viên giảng dạy, vị trí các tiết học đã được bố trí có thể được giữ nguyên và giáo viên của các tiết học được cập nhật theo thông tin phân công mới. Nếu giáo viên mới đã được bố trí giảng dạy lớp khác tại cùng thời điểm, hệ thống thông báo các tiết bị xung đột và xác định thời khóa biểu là cần điều chỉnh.

Đối với trường hợp thay đổi định mức số tiết của môn học, hệ thống xác định số tiết đang thiếu hoặc dư so với định mức mới. Khi định mức tăng, người dùng cần bổ sung các tiết còn thiếu. Khi định mức giảm, người dùng cần loại bỏ các tiết dư. Người dùng có thể điều chỉnh thời khóa biểu thủ công hoặc thực hiện lại chức năng sắp xếp tự động.

Nếu việc thay đổi dữ liệu làm phát sinh vi phạm các ràng buộc bắt buộc, hệ thống vẫn cho phép lưu dữ liệu mới nhưng thời khóa biểu được xác định là chưa hợp lệ và cần được kiểm tra lại. Các thay đổi có thể được áp dụng từ một tuần do người dùng lựa chọn, trong khi thời khóa biểu của các tuần trước đó vẫn được giữ nguyên.

Sau khi hoàn tất việc xây dựng và kiểm tra thời khóa biểu, người phụ trách thực hiện công bố để đưa thời khóa biểu vào sử dụng. Trước khi công bố, hệ thống kiểm tra tính đầy đủ của dữ liệu phân công giảng dạy, số tiết của từng môn học và toàn bộ các ràng buộc bắt buộc.

Nếu còn tồn tại vi phạm, hệ thống từ chối công bố và hiển thị danh sách các vấn đề cần xử lý. Người dùng có thể dựa trên thông tin được cung cấp để điều chỉnh thủ công hoặc thực hiện lại chức năng sắp xếp tự động.

Đối với thời khóa biểu hợp lệ, người phụ trách có thể thực hiện công bố để cung cấp cho người xem. Khi có thay đổi cần áp dụng trong quá trình giảng dạy, người dùng lựa chọn một tuần chưa công bố làm thời điểm bắt đầu áp dụng thời khóa biểu mới. Các tuần trước thời điểm này tiếp tục giữ nguyên lịch đã có, còn các tuần chưa công bố từ thời điểm được chọn trở đi sử dụng thời khóa biểu đã được điều chỉnh.

Đối với tuần đã được công bố, thời khóa biểu được khóa và giữ nguyên nội dung đã công bố. Những thay đổi đối với dữ liệu đầu vào chỉ được áp dụng khi xây dựng hoặc kiểm tra các tuần chưa công bố. Khi cần điều chỉnh một tuần đã công bố, người phụ trách phải hủy công bố, thực hiện chỉnh sửa và công bố lại sau khi thời khóa biểu đáp ứng toàn bộ các ràng buộc bắt buộc.

Người dùng vãng lai không cần đăng nhập và chỉ được phép xem thời khóa biểu đã được công bố, không thể truy cập hoặc chỉnh sửa dữ liệu quản lý của hệ thống.
% Vào đầu mỗi năm học, nhà trường tiến hành xây dựng thời khóa biểu cho các lớp nhằm đảm bảo hoạt động giảng dạy diễn ra theo đúng chương trình giáo dục và kế hoạch năm học. Công việc này thường do Ban Giám hiệu hoặc cán bộ phụ trách chuyên môn thực hiện dựa trên chương trình giảng dạy, đội ngũ giáo viên và các quy định của nhà trường.

% Quá trình xây dựng thời khóa biểu bắt đầu từ việc xác định các lớp học sẽ được tổ chức trong năm học, các môn học cần giảng dạy và số tiết của từng môn theo chương trình giáo dục. Tiếp theo, nhà trường thực hiện phân công giảng dạy, xác định giáo viên phụ trách từng môn học của từng lớp.

% Đối với trường tiểu học, việc phân công giảng dạy có một số đặc thù riêng. Thông thường, mỗi lớp được phân công một giáo viên chủ nhiệm phụ trách quản lý lớp và giảng dạy phần lớn các môn học văn hóa như Toán, Tiếng Việt, Đạo đức, Tự nhiên và Xã hội (tùy theo chương trình và quy định của từng trường). Các môn học chuyên biệt như Tiếng Anh, Tin học, Âm nhạc, Mỹ thuật và Giáo dục thể chất thường do giáo viên bộ môn đảm nhiệm. Một giáo viên bộ môn có thể được phân công giảng dạy nhiều lớp khác nhau, trong khi giáo viên chủ nhiệm chủ yếu giảng dạy lớp mình phụ trách.

% Sau khi hoàn tất việc phân công giảng dạy, người phụ trách tiến hành sắp xếp các tiết học vào các khung thời gian trong tuần để xây dựng thời khóa biểu cho từng lớp. Trong quá trình này cần đảm bảo nhiều yêu cầu như: mỗi giáo viên không được giảng dạy nhiều lớp trong cùng một thời điểm; mỗi lớp chỉ học một môn trong một tiết học; các môn học được bố trí đủ số tiết theo quy định của chương trình; các tiết học trong cùng một buổi được sắp xếp liên tục, không xuất hiện khoảng trống giữa các tiết học; đồng thời đáp ứng các ràng buộc do nhà trường quy định.

% Sau khi xây dựng được thời khóa biểu đáp ứng các yêu cầu bắt buộc, người phụ trách chuyên môn thường tiếp tục rà soát và điều chỉnh nhằm nâng cao chất lượng thời khóa biểu, bảo đảm cách phân bố các tiết học phù hợp với yêu cầu quản lý của nhà trường.

% Cuối cùng, thời khóa biểu được rà soát lần cuối để phát hiện và xử lý các xung đột hoặc những điểm chưa phù hợp trước khi chính thức ban hành và đưa vào áp dụng trong hoạt động giảng dạy.
% \section{Phân tích chức năng}
% \label{section:2.2}

% Hệ thống EduSchedule hỗ trợ ban giám hiệu quản lý và xếp thời khoá biểu cho toàn trường tiểu học. Sau đây là mô tả chi tiết các nghiệp vụ chính trong hệ thống.

% \textbf{Quản lý giáo viên}

% \textit{Nghiệp vụ}: Lưu trữ và theo dõi thông tin giáo viên trong nhà trường.

% \textit{Nghiệp vụ chi tiết}:
% \begin{itemize}[noitemsep]
%     \item Thêm, cập nhật hoặc xoá thông tin giáo viên, bao gồm họ tên, loại giáo viên (chủ nhiệm hoặc bộ môn) và định mức số tiết dạy tối đa mỗi tuần.
%     \item Đánh dấu giáo viên không còn hoạt động (inactive) thay vì xoá, nhằm bảo toàn dữ liệu phân công và thời khoá biểu các kỳ trước; giáo viên inactive không xuất hiện trong các lựa chọn phân công mới.
% \end{itemize}

% \textbf{Quản lý lớp học}

% \textit{Nghiệp vụ}: Lưu trữ danh sách lớp học theo từng năm học.

% \textit{Nghiệp vụ chi tiết}:
% \begin{itemize}[noitemsep]
%     \item Thêm, cập nhật hoặc xoá thông tin lớp học, bao gồm tên lớp và khối (khối 1 đến 5); mỗi lớp thuộc về một năm học cụ thể.
%     \item Gán giáo viên chủ nhiệm cho lớp; mỗi lớp có đúng một giáo viên chủ nhiệm trong một năm học.
% \end{itemize}

% \textbf{Quản lý môn học}

% \textit{Nghiệp vụ}: Lưu trữ danh mục môn học và số tiết yêu cầu theo từng khối.

% \textit{Nghiệp vụ chi tiết}:
% \begin{itemize}[noitemsep]
%     \item Thêm, cập nhật hoặc xoá thông tin môn học, bao gồm tên môn và số tiết yêu cầu mỗi tuần theo từng khối (khối 1 đến 5).
%     \item Danh mục môn học là cơ sở để thực hiện phân công giảng dạy và kiểm tra số tiết khi xếp lịch.
% \end{itemize}

% \textbf{Quản lý phòng chức năng}

% \textit{Nghiệp vụ}: Lưu trữ danh sách phòng chức năng và liên kết với môn học tương ứng.

% \textit{Nghiệp vụ chi tiết}:
% \begin{itemize}[noitemsep]
%     \item Thêm, cập nhật hoặc xoá phòng chức năng (phòng máy tính, phòng thể dục\ldots); mỗi phòng gắn với đúng một môn học và có thông tin số lượng phòng hiện có.
%     \item Khi xếp thời khoá biểu, hệ thống kiểm tra số lớp đang dùng phòng cùng một buổi và tiết không vượt quá số lượng phòng; nếu vi phạm, hệ thống thông báo ngay lập tức để người dùng điều chỉnh.
% \end{itemize}

% \textbf{Quản lý năm học}

% \textit{Nghiệp vụ}: Tạo và quản lý năm học, làm khung tổ chức cho phân công giảng dạy và thời khoá biểu.

% \textit{Nghiệp vụ chi tiết}:
% \begin{itemize}[noitemsep]
%     \item Tạo năm học mới bằng cách chọn năm bắt đầu và nhập ngày bắt đầu của học kì 1 và học kì 2; hệ thống tự động tính ngày cho từng tuần học (HK1: 18 tuần, HK2: 17 tuần theo quy định Bộ GD\&ĐT).
%     \item Khi tạo năm học, hệ thống tự động khởi tạo hai thời khoá biểu (học kì 1 và học kì 2) cùng danh sách tuần học tương ứng.
%     \item Cho phép điều chỉnh ngày bắt đầu của từng tuần học; thay đổi tuần N tự động cập nhật cascade cho các tuần N+1 trở đi.
% \end{itemize}

% \textbf{Phân công giảng dạy}

% \textit{Nghiệp vụ}: Xác định giáo viên phụ trách từng môn học cho từng lớp trong năm học.

% \textit{Nghiệp vụ chi tiết}:
% \begin{itemize}[noitemsep]
%     \item Gán giáo viên chủ nhiệm cho từng lớp; mỗi lớp phải có đúng một giáo viên chủ nhiệm trước khi tiến hành xếp thời khoá biểu.
%     \item Phân công giáo viên phụ trách từng môn học cho từng lớp; cả giáo viên chủ nhiệm lẫn giáo viên bộ môn đều có thể được phân công dạy bất kỳ môn học nào.
%     \item Tổng số tiết của mỗi giáo viên bằng tổng số tiết của tất cả các môn được phân công cho giáo viên đó; hệ thống hiển thị tổng này so với định mức và cảnh báo khi vượt định mức nhưng không chặn lưu.
% \end{itemize}

% \textbf{Xếp thời khoá biểu}

% \textit{Nghiệp vụ}: Phân bổ các tiết học vào khung giờ cụ thể trong tuần, hỗ trợ hai phương thức: điều chỉnh thủ công trên giao diện lưới và xếp tự động bằng thuật toán.

% \textit{Nghiệp vụ chi tiết}:
% \begin{itemize}[noitemsep]
%     \item \textit{Xếp thủ công}: người dùng nhấn vào ô trống trên lưới để chọn môn cần xếp; hệ thống hiển thị danh sách môn còn lại của lớp kèm tên giáo viên và số tiết còn thiếu trong tuần; môn đã đủ tiết bị vô hiệu hoá.
%     \item \textit{Xếp tự động}: hệ thống sinh lịch cho các môn chưa được xếp, giữ nguyên các tiết đã xếp thủ công trước đó; người dùng xem xét kết quả và quyết định có áp dụng hay không.
%     \item Một giáo viên không thể dạy hai lớp khác nhau trong cùng một buổi và tiết; hệ thống phát hiện và thông báo xung đột ngay lập tức.
%     \item Số lớp sử dụng phòng chức năng cho cùng một môn trong cùng một buổi và tiết không được vượt quá số lượng phòng hiện có; hệ thống thông báo vi phạm để người dùng điều chỉnh.
%     \item Sau khi hoàn tất một tuần, người dùng có thể lưu riêng tuần đó hoặc áp dụng lịch đó cho tất cả các tuần tiếp theo trong kì.
% \end{itemize}

% \textbf{Xem và xuất thời khoá biểu}

% \textit{Nghiệp vụ}: Tra cứu thời khoá biểu và xuất file phục vụ in ấn.

% \textit{Nghiệp vụ chi tiết}:
% \begin{itemize}[noitemsep]
%     \item Xem thời khoá biểu theo lớp, theo khối hoặc theo giáo viên; chọn tuần học cụ thể để xem.
%     \item Các ô đang có xung đột được đánh dấu trực quan để người dùng nhận biết và điều chỉnh.
%     \item Xuất thời khoá biểu ra file Excel để in ấn và phân phát cho toàn trường.
% \end{itemize}

% Biểu đồ hoạt động minh hoạ toàn bộ quy trình xếp thời khoá biểu từ khâu chuẩn bị đến khi hoàn thiện được thể hiện trong Hình~\ref{fig:activity_xep_tkb_admin}.

% \begin{figure}[H]
% 	\centering
% 	\includegraphics[width=0.9\textwidth]{Hinhve/Quy_trinh_nghiep_vu/Nghiệp vụ xếp TKB của admin.jpg}
% 	\caption{Biểu đồ hoạt động quy trình xếp thời khoá biểu}
% 	\label{fig:activity_xep_tkb_admin}
% \end{figure}

% Tạm ẩn phần "Phân tích đầu vào của hệ thống" để tinh gọn Chương 2.
% Giữ lại nội dung để tham khảo; xóa hai lệnh \iffalse và \fi nếu muốn
% đưa phần này trở lại báo cáo.
\iffalse
\section{Phân tích đầu vào của hệ thống}
\label{section:2.3}

\subsection{Năm học (SchoolYear)}
\label{section:2.3.1}
Năm học là đơn vị tổ chức cao nhất, làm ngữ cảnh cho toàn bộ dữ liệu lớp học, phân công và thời khoá biểu. Các thông tin chính bao gồm:
\begin{itemize}[noitemsep]
	\item Năm bắt đầu (ví dụ: 2025 → năm học 2025--2026)
	\item Ngày bắt đầu học kì 1 và học kì 2
\end{itemize}
Mỗi năm học chứa hai thời khoá biểu (học kì 1 và học kì 2) được khởi tạo tự động khi tạo năm học.

\subsection{Thông tin giáo viên (Teacher)}
\label{section:2.3.2}
Dữ liệu giáo viên được nhập thủ công qua giao diện quản lý. Các thông tin chính bao gồm:
\begin{itemize}[noitemsep]
	\item Họ và tên
	\item Loại giáo viên: chủ nhiệm hoặc bộ môn
	\item Định mức số tiết tối đa mỗi tuần
	\item Trạng thái hoạt động (active/inactive)
\end{itemize}

\subsection{Thông tin lớp học (SchoolClass)}
\label{section:2.3.3}
Lớp học gắn với một năm học cụ thể. Các thông tin chính bao gồm:
\begin{itemize}[noitemsep]
	\item Tên lớp (ví dụ: 4A, 4B)
	\item Khối (từ 1 đến 5)
	\item Giáo viên chủ nhiệm: liên kết tới Teacher
	\item Năm học: liên kết tới SchoolYear
\end{itemize}

\subsection{Thông tin môn học (Subject)}
\label{section:2.3.4}
Danh mục môn học dùng chung, không phân biệt năm học. Các thông tin chính bao gồm:
\begin{itemize}[noitemsep]
	\item Tên môn học (ví dụ: Toán, Tiếng Việt, Tiếng Anh)
	\item Số tiết theo khối (ví dụ: khối 1: 1 tiết/tuần, khối 4: 4 tiết/tuần)
\end{itemize}

\subsection{Phòng chức năng (SpecialRoom)}
\label{section:2.3.5}
Phòng chức năng được quản lý toàn cục, không phân biệt năm học. Các thông tin chính bao gồm:
\begin{itemize}[noitemsep]
	\item Tên phòng (ví dụ: Phòng máy tính 1, Phòng thể dục)
	\item Môn học liên kết: mỗi phòng gắn với đúng một môn
	\item Số lượng phòng hiện có
\end{itemize}

\subsection{Phân công giảng dạy (Assignment)}
\label{section:2.3.6}
Phân công xác định giáo viên phụ trách một môn học cho một lớp trong năm học. Các thông tin chính bao gồm:
\begin{itemize}[noitemsep]
	\item Lớp học, môn học, giáo viên phụ trách
	\item Số tiết mỗi tuần theo phân công
\end{itemize}
Ràng buộc: mỗi cặp (lớp, môn, giáo viên) là duy nhất; phân công ngầm thuộc năm học thông qua lớp học.

\subsection{Tuần học (Week) và Tiết học (Slot)}
\label{section:2.3.7}
Mỗi thời khoá biểu gồm nhiều tuần học. Tuần học lưu ngày bắt đầu (thứ Hai) và liên kết tới thời khoá biểu. Mỗi tiết học (Slot) thuộc một tuần và lưu các thông tin:
\begin{itemize}[noitemsep]
	\item Phân công tương ứng (liên kết tới Assignment — bao gồm lớp, môn, giáo viên)
	\item Thứ trong tuần, buổi học (sáng/chiều), tiết trong buổi
	\item Phòng chức năng sử dụng (nếu có)
\end{itemize}
\fi

% Tạm ẩn phần "Đầu ra của hệ thống" vì nội dung bị trùng với các phần
% mô tả nghiệp vụ, đặc tả use case và yêu cầu phi chức năng. Giữ lại nội dung
% để tham khảo; xóa hai lệnh \iffalse và \fi nếu muốn đưa phần này trở lại.
\iffalse
\section{Đầu ra của hệ thống}
\label{section:2.4}

\subsection{Thời khóa biểu dạng lưới (Timetable Grid)}
\label{subsection:2.3.1}
Hiển thị lịch dạy dưới dạng bảng lưới với hàng là tiết học (tiết 1 đến tiết 7) và cột là thứ trong tuần (thứ Hai đến thứ Sáu). Hệ thống hỗ trợ ba chế độ xem: theo lớp, theo khối và theo giáo viên. Mỗi ô trong lưới hiển thị tên môn học và tên giáo viên bộ môn phụ trách (nếu có).

\subsection{Hiển thị xung đột (Conflict Display)}
Hệ thống hiển thị xung đột thời khóa biểu thông qua hai trạng thái:
\begin{itemize}[noitemsep]
	\item Khi xem thời khóa biểu, các ô đang có xung đột được tô màu đỏ giúp người dùng nhận biết ngay vị trí cần điều chỉnh;
	\item Khi đang thao tác xếp tiết mà tiết vừa chọn gây ra xung đột, ô đó lập tức chuyển sang màu đỏ và hệ thống hiển thị thông báo mô tả chi tiết thông tin.
\end{itemize}

\subsection{File xuất thời khóa biểu (Export File)}
\label{subsection:2.4.4}
Hệ thống cho phép xuất thời khóa biểu ra file Excel theo mẫu chuẩn, phân chia theo thứ và tiết, tên giáo viên bộ môn được hiển thị rõ ràng. Tính năng hỗ trợ xuất theo lớp hoặc theo giáo viên tùy nhu cầu sử dụng.
\fi

\section{Tổng quan chức năng}
\label{section:2.3}
% Phần này tóm tắt các chức năng của phần mềm ở mức cao. Đặc tả chi tiết được trình bày trong phần tiếp theo.

\subsection{Biểu đồ ca sử dụng tổng quát}
\label{subsection:2.3.1}
% Sinh viên vẽ biểu đồ use case tổng quan và giải thích các tác nhân tham gia là gì, nêu vai trò của từng tác nhân, và mô tả ngắn gọn các use case chính.
\begin{figure}[H]
	\centering
	\includegraphics[width=0.9\textwidth]{Hinhve/UsecaseDiagram/Sơ đồ usecase tổng quát EduSchedule.jpg}
	\caption{Biểu đồ Use Case hệ thống}
    
	\label{fig:usecase}
\end{figure}

Hình~\ref{fig:usecase} minh họa biểu đồ ca sử dụng tổng quát của hệ thống
EduSchedule. Hệ thống có hai tác nhân chính là Người xếp thời khóa biểu và
Người dùng vãng lai. Mỗi tác nhân tương tác với hệ thống theo phạm vi chức
năng và quyền truy cập riêng.

Người xếp thời khóa biểu là tác nhân trực tiếp thực hiện các nghiệp vụ quản lý
và xây dựng thời khóa biểu của nhà trường. Sau khi đăng nhập, tác nhân này có
thể quản lý năm học, giáo viên, lớp học, môn học và phòng chức năng; thực hiện
phân công giảng dạy; xếp thời khóa biểu; công khai thời khóa biểu và xuất kết
quả ra tệp Excel. Trong đó, phân công giảng dạy và xếp thời khóa biểu là hai
nhóm chức năng trọng tâm. Các dữ liệu được quản lý là cơ sở để hệ thống kiểm
tra ràng buộc và hỗ trợ người dùng xây dựng thời khóa biểu phù hợp.

Người dùng vãng lai là giáo viên hoặc cá nhân có nhu cầu tra cứu thời khóa
biểu đã được công khai. Tác nhân này không cần đăng nhập và không có quyền
thay đổi dữ liệu của hệ thống. Việc tách riêng chức năng tra cứu giúp người dùng
tiếp cận thời khóa biểu thuận tiện, đồng thời bảo đảm các chức năng quản lý và
xếp lịch chỉ được thực hiện bởi người có trách nhiệm.

\subsection{Biểu đồ ca sử dụng phân rã phân công giảng dạy}
\begin{figure}[H]
	\centering
	\includegraphics[width=0.9\textwidth]{Hinhve/UsecaseDiagram/Sơ đồ phân rã Phân công giảng dạy.jpg}
	\caption{Biểu đồ ca sử dụng phân rã phân công giảng dạy}
	\label{fig:usecase_phan_cong_giang_day}
\end{figure}

Hình~\ref{fig:usecase_phan_cong_giang_day} mô tả các chức năng thuộc nghiệp
vụ phân công giảng dạy do Người xếp thời khóa biểu thực hiện. Nghiệp vụ này
được chia thành hai trường hợp sử dụng là Phân công giáo viên chủ nhiệm và
Phân công giáo viên bộ môn. Với chức năng Phân công giáo viên chủ nhiệm,
người dùng lựa chọn giáo viên phụ trách một lớp trong năm học. Với chức năng
Phân công giáo viên bộ môn, người dùng xác định giáo viên giảng dạy từng môn
học cho từng lớp. Kết quả phân công là dữ liệu đầu vào quan trọng để xác định
giáo viên phụ trách mỗi tiết học và kiểm tra xung đột giáo viên trong quá trình
xếp thời khóa biểu.

\subsection{Biểu đồ ca sử dụng phân rã xếp thời khoá biểu}
\begin{figure}[H]
	\centering
	\includegraphics[width=0.9\textwidth]{Hinhve/UsecaseDiagram/Sơ đồ phân rã Xếp thời khoá biểu.jpg}
	\caption{Biểu đồ ca sử dụng phân rã xếp thời khóa biểu}
	\label{fig:usecase_xep_thoi_khoa_bieu}
\end{figure}

Hình~\ref{fig:usecase_xep_thoi_khoa_bieu} mô tả hai phương thức xếp thời
khóa biểu mà Người xếp thời khóa biểu có thể lựa chọn. Với phương thức Xếp
thời khóa biểu thủ công, người dùng chủ động lựa chọn môn học và bố trí từng
tiết vào các vị trí trên lưới thời khóa biểu; hệ thống kiểm tra và thông báo khi
phát sinh vi phạm ràng buộc. Với phương thức Xếp thời khóa biểu tự động, hệ
thống sử dụng dữ liệu phân công giảng dạy và các ràng buộc đã thiết lập để tạo
phương án cho những tiết học chưa được sắp xếp, đồng thời giữ nguyên các tiết
đã được người dùng bố trí trước đó. Kết quả tự động được hiển thị để người dùng
kiểm tra và xác nhận trước khi áp dụng. Hai phương thức có thể được sử dụng
kết hợp nhằm vừa hỗ trợ tự động hóa quá trình lập lịch, vừa bảo đảm người dùng
có thể chủ động điều chỉnh thời khóa biểu theo nhu cầu thực tế.


\section{Đặc tả chức năng}
\label{section:2.4}

Phần này trình bày đặc tả các chức năng chính của hệ thống EduSchedule, bao gồm Đăng nhập, Phân công giảng dạy, Xếp thời khoá biểu và Tra cứu thời khoá biểu. Mỗi đặc tả gồm mã và tên ca sử dụng, tác nhân, tiền điều kiện, luồng sự kiện chính, luồng sự kiện thay thế và hậu điều kiện. Đặc tả các chức năng còn lại của hệ thống được trình bày tại Phụ lục~A.
\subsection{Đặc tả ca sử dụng: Đăng nhập}

Bảng~\ref{tab:uc_dang_nhap} đặc tả ca sử dụng đăng nhập, đối tượng áp dụng là Người xếp thời khóa biểu, với mục đích xác thực để truy cập các chức năng quản lý của hệ thống.

\begin{longtable}{|p{2.8cm}|p{1cm}|p{2.8cm}|p{6.3cm}|}
\caption{Đặc tả ca sử dụng Đăng nhập}
\label{tab:uc_dang_nhap} \\
\hline
\multicolumn{1}{|l|}{\textbf{Mã use case}} &
\multicolumn{1}{c|}{UC01} &
\multicolumn{1}{c|}{\textbf{Tên use case}} &
\multicolumn{1}{c|}{Đăng nhập} \\
\hline
\multicolumn{1}{|l|}{\textbf{Tác nhân}} &
\multicolumn{3}{l|}{Người xếp thời khóa biểu} \\
\hline
\multicolumn{1}{|l|}{\textbf{Tiền điều kiện}} &
\multicolumn{3}{l|}{Không} \\
\hline
\multirow{1}{2.8cm}{Luồng sự kiện chính} &
\multicolumn{3}{c|}{%
\begin{tabular}{|p{0.8cm}|p{2.8cm}|p{6.2cm}|}
\hline
\textbf{STT} & \textbf{Thực hiện bởi} & \textbf{Hành động} \\
\hline
1. & Người dùng & Truy cập trang đăng nhập của hệ thống. \\
\hline
2. & Hệ thống & Hiển thị giao diện đăng nhập. \\
\hline
3. & Người dùng & Nhập tên đăng nhập, mật khẩu và nhấn nút ``Đăng nhập''. \\
\hline
4. & Hệ thống & Đăng nhập thành công và chuyển hướng tới trang chủ. \\
\hline
\end{tabular}} \\
\hline
\multirow{1}{2.8cm}{Luồng sự kiện thay thế} &
\multicolumn{3}{c|}{%
\begin{tabular}{|p{0.8cm}|p{2.8cm}|p{6.2cm}|}
\hline
\textbf{STT} & \textbf{Thực hiện bởi} & \textbf{Hành động} \\
\hline
4a. & Hệ thống & Thông báo lỗi nếu tên đăng nhập hoặc mật khẩu không chính xác. \\
\hline
4b. & Hệ thống & Thông báo lỗi nếu người dùng chưa nhập đầy đủ tên đăng nhập hoặc mật khẩu. \\
\hline
\end{tabular}} \\
\hline
\multicolumn{1}{|l|}{\textbf{Hậu điều kiện}} &
\multicolumn{3}{l|}{Không} \\
\hline
\end{longtable}

\subsection{Đặc tả ca sử dụng: Phân công giáo viên chủ nhiệm}
Bảng~\ref{tab:uc_phan_cong_chu_nhiem} đặc tả ca sử dụng phân công giáo viên chủ nhiệm, đối tượng áp dụng là Người xếp thời khóa biểu, với mục đích xác định giáo viên chủ nhiệm cho từng lớp trong năm học cụ thể.

\begin{longtable}{|p{2.8cm}|p{1cm}|p{2.8cm}|p{6.3cm}|}
\caption{Đặc tả ca sử dụng Phân công giáo viên chủ nhiệm}
\label{tab:uc_phan_cong_chu_nhiem} \\
\hline
\multicolumn{1}{|l|}{\textbf{Mã use case}} & \multicolumn{1}{c|}{UC02} & \multicolumn{1}{c|}{\textbf{Tên use case}} & \multicolumn{1}{c|}{Phân công giáo viên chủ nhiệm} \\
\hline
\multicolumn{1}{|l|}{\textbf{Tác nhân}} & \multicolumn{3}{l|}{Người xếp thời khoá biểu} \\
\hline
\multicolumn{1}{|l|}{\textbf{Tiền điều kiện}} & \multicolumn{3}{p{10.5cm}|}{
\begin{enumerate}[leftmargin=*, noitemsep, topsep=0pt, partopsep=0pt, parsep=0pt]
    \item Người dùng đã đăng nhập thành công.
    \item Đã tạo năm học; danh sách giáo viên và danh sách lớp học.
\end{enumerate}} \\
\hline
\multirow{1}{2.8cm}{Luồng sự kiện chính} &
\multicolumn{3}{c|}{%
\begin{tabular}{|p{0.8cm}|p{2.8cm}|p{6.2cm}|}
\hline
\textbf{STT} & \textbf{Thực hiện bởi} & \textbf{Hành động} \\
\hline
1. & Người dùng & Vào mục ``Phân công giảng dạy'', chọn tab ``Phân công Chủ nhiệm''. \\
\hline
2. & Người dùng & Chọn một lớp và chọn giáo viên chủ nhiệm từ danh sách; danh sách hiển thị tất cả giáo viên, trừ những giáo viên hiện đã là GVCN của lớp khác trong năm học. \\
\hline
3. & Hệ thống & Lưu thông tin và cập nhật giao diện. \\
\hline
\end{tabular}} \\
\hline
\multirow{1}{2.8cm}{Luồng sự kiện thay thế} &
\multicolumn{3}{c|}{%
\begin{tabular}{|p{0.8cm}|p{2.8cm}|p{6.2cm}|}
\hline
\textbf{STT} & \textbf{Thực hiện bởi} & \textbf{Hành động} \\
\hline
2a. & Người dùng & Thay vì chọn giáo viên mới, chọn xoá giáo viên chủ nhiệm hiện tại của lớp (chọn giá trị trống). \\
\hline
2b. & Hệ thống & Hiển thị hộp thoại xác nhận xoá GVCN của lớp; nếu người dùng xác nhận, hệ thống xoá giáo viên chủ nhiệm khỏi lớp và cập nhật giao diện. \\
\hline
\end{tabular}} \\
\hline
\multicolumn{1}{|l|}{\textbf{Hậu điều kiện}} & \multicolumn{3}{p{10.5cm}|}{Giáo viên chủ nhiệm được gán cho lớp; dữ liệu được lưu vào cơ sở dữ liệu.} \\
\hline
\end{longtable}

\subsection{Đặc tả ca sử dụng: Phân công giáo viên bộ môn}
Bảng~\ref{tab:uc_phan_cong_bo_mon} đặc tả ca sử dụng phân công giáo viên bộ môn, đối tượng áp dụng là Người xếp thời khóa biểu, với mục đích xác định giáo viên phụ trách từng môn học cho từng lớp trong năm học cụ thể.

\begin{longtable}{|p{2.8cm}|p{1cm}|p{2.8cm}|p{6.3cm}|}
\caption{Đặc tả ca sử dụng Phân công giáo viên bộ môn}
\label{tab:uc_phan_cong_bo_mon} \\
\hline
\multicolumn{1}{|l|}{\textbf{Mã use case}} & \multicolumn{1}{c|}{UC03} & \multicolumn{1}{c|}{\textbf{Tên use case}} & \multicolumn{1}{c|}{Phân công giáo viên bộ môn} \\
\hline
\multicolumn{1}{|l|}{\textbf{Tác nhân}} & \multicolumn{3}{l|}{Người xếp thời khoá biểu} \\
\hline
\multicolumn{1}{|l|}{\textbf{Tiền điều kiện}} & \multicolumn{3}{p{10.5cm}|}{
\begin{enumerate}[leftmargin=*, noitemsep, topsep=0pt, partopsep=0pt, parsep=0pt]
    \item Người dùng đã đăng nhập thành công.
    \item Đã tạo năm học; danh sách giáo viên, lớp học và môn học đã được thiết lập.
\end{enumerate}} \\
\hline
\multirow{1}{2.8cm}{Luồng sự kiện chính} &
\multicolumn{3}{c|}{%
\begin{tabular}{|p{0.8cm}|p{2.8cm}|p{6.2cm}|}
\hline
\textbf{STT} & \textbf{Thực hiện bởi} & \textbf{Hành động} \\
\hline
1. & Người dùng & Vào mục ``Phân công giảng dạy'', chọn tab ``Phân công Chuyên môn''. \\
\hline
2. & Người dùng & Chọn hiển thị theo giáo viên hoặc theo lớp. \\
\hline
3. & Hệ thống & Hiển thị danh sách giáo viên hoặc lớp học tương ứng; \\
\hline
4. & Người dùng & Chọn một môn học của một lớp và chọn giáo viên phụ trách từ danh sách. \\
\hline
5. & Hệ thống & Lưu phân công và cập nhật giao diện. \\
\hline
6. & Hệ thống & Cập nhật tổng số tiết đã phân công của giáo viên đó so với định mức. \\
\hline
\end{tabular}} \\
\hline
\multirow{1}{2.8cm}{Luồng sự kiện thay thế} &
\multicolumn{3}{c|}{%
\begin{tabular}{|p{0.8cm}|p{2.8cm}|p{6.2cm}|}
\hline
\textbf{STT} & \textbf{Thực hiện bởi} & \textbf{Hành động} \\
\hline
4a. & Hệ thống & Nếu môn học của lớp đã được phân công cho một giáo viên khác trước đó, hệ thống tự động thay thế bằng giáo viên mới được chọn ngay khi lưu, không hiển thị xác nhận trước khi ghi đè. \\
\hline
6a. & Hệ thống & Khi tổng số tiết đã phân công của giáo viên vượt định mức, số tiết hiển thị màu đỏ trong bảng ``Theo giáo viên''; hệ thống không chặn lưu. \\
\hline
\end{tabular}} \\
\hline
\multicolumn{1}{|l|}{\textbf{Hậu điều kiện}} & \multicolumn{3}{p{10.5cm}|}{Bản ghi phân công được lưu và tổng số tiết phân công của giáo viên được cập nhật.} \\
\hline
\end{longtable}

\subsection{Đặc tả ca sử dụng: Xếp và điều chỉnh thời khoá biểu thủ công}

Bảng~\ref{tab:uc_xep_tkb_thu_cong} đặc tả ca sử dụng xếp và điều chỉnh thời khoá biểu thủ công, đối tượng áp dụng là Người xếp thời khoá biểu. Ca sử dụng cho phép người dùng tiếp tục hoàn thiện thời khoá biểu đã có một số tiết, bố trí môn học vào từng ô trên lưới và điều chỉnh các tiết không còn phù hợp với thông tin phân công giảng dạy hiện tại.

\begin{longtable}{|p{2.8cm}|p{1cm}|p{2.8cm}|p{6.3cm}|}
\caption{Đặc tả ca sử dụng Xếp và điều chỉnh thời khoá biểu thủ công}
\label{tab:uc_xep_tkb_thu_cong} \\
\hline
\multicolumn{1}{|l|}{\textbf{Mã use case}} & \multicolumn{1}{c|}{UC04} & \multicolumn{1}{c|}{\textbf{Tên use case}} & \multicolumn{1}{c|}{Xếp và điều chỉnh thời khoá biểu thủ công} \\
\hline
\multicolumn{1}{|l|}{\textbf{Tác nhân}} & \multicolumn{3}{l|}{Người xếp thời khoá biểu} \\
\hline
\multicolumn{1}{|l|}{\textbf{Tiền điều kiện}} & \multicolumn{3}{p{10.5cm}|}{
\begin{enumerate}[noitemsep, topsep=2pt]
    \item Người dùng đã đăng nhập thành công.
    \item Đã tạo năm học và thời khoá biểu tương ứng.
    \item Các lớp học, môn học và phân công giảng dạy đã được thiết lập đầy đủ.
\end{enumerate}} \\
\hline
\multirow{1}{2.8cm}{Luồng sự kiện chính} &
\multicolumn{3}{c|}{%
\begin{tabular}{|p{0.8cm}|p{2.8cm}|p{6.2cm}|}
\hline
\textbf{STT} & \textbf{Thực hiện bởi} & \textbf{Hành động} \\
\hline
1. & Người dùng & Truy cập trang ``Thời khoá biểu'' và chọn năm học, học kì và tuần cần xếp. \\
\hline
2. & Hệ thống & Tải các tiết đã được xếp và kiểm tra lại giáo viên của từng tiết theo thông tin phân công giảng dạy hiện tại. \\
\hline
3. & Người dùng & Chọn một ô trống hoặc một ô đã có tiết học trên lưới (xác định thứ, buổi và tiết). \\
\hline
4. & Hệ thống & Hiển thị danh sách môn học tương ứng và số tiết chưa xếp của lớp đó. \\
\hline
5. & Người dùng & Chọn môn học cần bố trí hoặc thay thế trong ô. \\
\hline
6. & Hệ thống & Kiểm tra xung đột giáo viên, phòng chức năng và các ràng buộc liên quan tại vị trí được chọn. \\
\hline
7. & Hệ thống & Cập nhật tiết học trên giao diện lưới ở trạng thái chưa lưu. \\
\hline
8. & Người dùng & Tiếp tục xếp hoặc điều chỉnh các tiết còn lại cho đến khi thời khoá biểu được hoàn thiện. \\
\hline
9. & Người dùng & Nhấn ``Lưu tuần N (ngày)'' để lưu tuần đang chỉnh sửa, hoặc ``Áp dụng từ tuần N (ngày) trở đi'' để áp dụng thời khoá biểu cho tuần đó và các tuần tiếp theo. \\
\hline
10. & Hệ thống & Lưu các thay đổi theo phạm vi được người dùng lựa chọn và cập nhật giao diện lưới. \\
\hline
\end{tabular}} \\
\hline
\multirow{1}{2.8cm}{Luồng sự kiện thay thế} &
\multicolumn{3}{c|}{%
\begin{tabular}{|p{0.8cm}|p{2.8cm}|p{6.2cm}|}
\hline
\textbf{STT} & \textbf{Thực hiện bởi} & \textbf{Hành động} \\
\hline
2a. & Hệ thống & Nếu việc thay đổi phân công làm cho giáo viên hiện tại của một tiết học bị trùng lịch với lớp khác, hệ thống làm nổi bật các ô bị xung đột và hiển thị thông báo nêu rõ giáo viên, lớp học và thời điểm xảy ra xung đột để người dùng điều chỉnh. \\
\hline
3a. & Hệ thống & Nếu ô đã có tiết học, hệ thống cho phép người dùng chọn môn học khác để thay thế hoặc loại bỏ tiết học; các ràng buộc được kiểm tra lại sau khi điều chỉnh. \\
\hline
4a. & Hệ thống & Nếu môn học chưa được phân công giáo viên, hệ thống hiển thị thông báo chưa được phân công và không cho phép lựa chọn môn học đó. \\
\hline
4b. & Hệ thống & Nếu môn học đã được xếp đủ số tiết, hệ thống hiển thị trạng thái đã đủ và không cho phép người dùng lựa chọn môn học đó. \\
\hline
6a. & Hệ thống & Nếu giáo viên đang dạy lớp khác vào tiết này, hệ thống hiển thị thông báo xung đột và cho phép người dùng lựa chọn loại bỏ tiết học xung đột ở lớp khác trước khi lưu tiết mới. \\
\hline
6b. & Hệ thống & Nếu môn học cần phòng chức năng nhưng không có phòng trống, hiển thị thông báo xung đột phòng và không cho phép xếp môn học vào ô đó. \\
\end{tabular}} \\
\hline
\multicolumn{1}{|l|}{\textbf{Hậu điều kiện}} & \multicolumn{3}{p{10.5cm}|}{Các tiết học được lưu với thông tin thời khoá biểu, thứ, tiết và phân công tương ứng. Giao diện lưới phản ánh trạng thái mới nhất; các tiết còn xung đột với phân công hiện tại được làm nổi bật kèm thông báo lỗi để người dùng tiếp tục điều chỉnh.} \\
\hline
\end{longtable}

\subsection{Đặc tả ca sử dụng: Xếp thời khoá biểu tự động}

Bảng~\ref{tab:uc_xep_tkb_tu_dong} đặc tả ca sử dụng xếp thời khoá biểu tự động, đối tượng áp dụng là Người xếp thời khoá biểu, với mục đích tự động phân bổ các tiết học còn trống bằng thuật toán, giữ nguyên các tiết đã xếp thủ công.

\begin{longtable}{|p{2.8cm}|p{1cm}|p{2.8cm}|p{6.3cm}|}
\caption{Đặc tả ca sử dụng Xếp thời khoá biểu tự động}
\label{tab:uc_xep_tkb_tu_dong} \\
\hline
\multicolumn{1}{|l|}{\textbf{Mã use case}} & \multicolumn{1}{c|}{UC05} & \multicolumn{1}{c|}{\textbf{Tên use case}} & \multicolumn{1}{c|}{Xếp thời khoá biểu tự động} \\
\hline
\multicolumn{1}{|l|}{\textbf{Tác nhân}} & \multicolumn{3}{l|}{Người xếp thời khoá biểu} \\
\hline
\multicolumn{1}{|l|}{\textbf{Tiền điều kiện}} & \multicolumn{3}{p{10.5cm}|}{
\begin{enumerate}[noitemsep, topsep=2pt]
    \item Người dùng đã đăng nhập thành công.
\end{enumerate}} \\
\hline
\multirow{1}{2.8cm}{Luồng sự kiện chính} &
\multicolumn{3}{c|}{%
\begin{tabular}{|p{0.8cm}|p{2.8cm}|p{6.2cm}|}
\hline
\textbf{STT} & \textbf{Thực hiện bởi} & \textbf{Hành động} \\
\hline
1. & Người dùng & Truy cập trang ``Thời khoá biểu'' và chọn năm học, học kì, tuần cần xếp. \\
\hline
2. & Người dùng & Nhấn nút ``Tự động xếp TKB''. \\
\hline
3. & Hệ thống & Hiển thị màn hình chờ toàn trang trong khi thuật toán xử lý. \\
\hline
4. & Hệ thống & Chạy thuật toán xếp lịch tự động để phân bổ các tiết còn thiếu, giữ nguyên các tiết đã xếp trước đó. \\
\hline
5. & Hệ thống & Trả về danh sách tiết đã xếp ở trạng thái chưa lưu và danh sách tiết không thể xếp được (nếu có). \\
\hline
6. & Người dùng & Xem xét kết quả trên lưới thời khoá biểu. \\
\hline
7. & Người dùng & Nhấn ``Lưu tuần N (ngày)'' để lưu tuần đang chỉnh sửa, hoặc ``Áp dụng từ tuần N (ngày) trở đi'' để áp dụng kết quả cho tuần đó và các tuần tiếp theo. \\
\hline
8. & Hệ thống & Lưu kết quả xếp thời khoá biểu theo phạm vi người dùng đã lựa chọn và cập nhật lưới thời khoá biểu. \\
\hline
\end{tabular}} \\
\hline
\multirow{1}{2.8cm}{Luồng sự kiện thay thế} &
\multicolumn{3}{c|}{%
\begin{tabular}{|p{0.8cm}|p{2.8cm}|p{6.2cm}|}
\hline
\textbf{STT} & \textbf{Thực hiện bởi} & \textbf{Hành động} \\
\hline
2a. & Hệ thống & Nếu còn lớp hoặc môn học chưa có giáo viên phụ trách, hiển thị thông báo liệt kê danh sách và yêu cầu hoàn tất phân công trước; không gửi yêu cầu xếp tự động. \\
\hline
5a. & Hệ thống & Nếu có tiết không thể xếp được, vẫn hiển thị các tiết đã xếp ở trạng thái chưa lưu kèm thông báo cảnh báo số tiết đã xếp và số môn không xếp được; người dùng có thể lưu kết quả một phần và xử lý thủ công các tiết còn lại. \\
\hline
7a. & Người dùng & Nhấn nút đóng màn hình chỉnh sửa hoặc chuyển sang tuần khác khi còn thay đổi chưa lưu; hệ thống hiển thị hộp thoại xác nhận huỷ thay đổi. \\
\hline
7b. & Hệ thống & Nếu người dùng xác nhận huỷ, toàn bộ kết quả tự động bị loại bỏ, thời khoá biểu giữ nguyên trạng thái trước khi xếp; nếu chọn ở lại, kết quả chưa lưu vẫn được giữ nguyên trên giao diện. \\
\hline
\end{tabular}} \\
\hline
\multicolumn{1}{|l|}{\textbf{Hậu điều kiện}} & \multicolumn{3}{p{10.5cm}|}{Các tiết học được lưu vào cơ sở dữ liệu. Giao diện lưới cập nhật phản ánh thời khoá biểu mới. Các tiết không xếp được vẫn để trống và có thể tiếp tục xếp thủ công.} \\
\hline
\end{longtable}

\subsection{Đặc tả ca sử dụng: Tra cứu thời khoá biểu}

Bảng~\ref{tab:uc_tra_cuu_tkb} đặc tả ca sử dụng tra cứu thời khoá biểu, đối
tượng áp dụng là Người dùng vãng lai. Khi Người xếp thời khoá biểu thực hiện
công khai một thời khoá biểu, hệ thống sinh một UID duy nhất và tạo URL công
khai tương ứng. URL này xác định thời khoá biểu của một học kì thuộc một năm
học cụ thể, cho phép Người dùng vãng lai truy cập và xem theo lớp, theo khối
hoặc theo giáo viên mà không cần đăng nhập.

\begin{longtable}{|p{2.8cm}|p{1cm}|p{2.8cm}|p{6.3cm}|}
\caption{Đặc tả ca sử dụng Tra cứu thời khoá biểu}
\label{tab:uc_tra_cuu_tkb} \\
\hline
\multicolumn{1}{|l|}{\textbf{Mã use case}} & \multicolumn{1}{c|}{UC06} & \multicolumn{1}{c|}{\textbf{Tên use case}} & \multicolumn{1}{c|}{Tra cứu thời khoá biểu} \\
\hline
\multicolumn{1}{|l|}{\textbf{Tác nhân}} & \multicolumn{3}{l|}{Người dùng vãng lai} \\
\hline
\multicolumn{1}{|l|}{\textbf{Tiền điều kiện}} & \multicolumn{3}{p{10.5cm}|}{Thời khoá biểu đã được Người xếp thời khoá biểu công khai; hệ thống đã sinh UID và URL công khai tương ứng.} \\
\hline
\multirow{1}{2.8cm}{\textbf{Luồng sự kiện chính}} &
\multicolumn{3}{c|}{%
\begin{tabular}{|p{0.8cm}|p{2.8cm}|p{6.2cm}|}
\hline
\textbf{STT} & \textbf{Thực hiện bởi} & \textbf{Hành động} \\
\hline
1. & Người dùng vãng lai & Mở URL công khai của thời khoá biểu có chứa UID tương ứng. \\
\hline
2. & Hệ thống & Tiếp nhận UID từ URL và tra cứu thời khoá biểu đã được công khai. \\
\hline
3. & Hệ thống & Hiển thị thời khoá biểu công khai dưới dạng lưới tương ứng với năm học và học kì cố định. \\
\hline
4. & Người dùng vãng lai & Chọn tuần học, chế độ xem và lớp, khối hoặc giáo viên cần tra cứu. \\
\hline
5. & Hệ thống & Hiển thị thời khoá biểu đã được công khai dưới dạng lưới theo lựa chọn; mỗi ô thể hiện môn học và giáo viên phụ trách. \\
\hline
\end{tabular}} \\
\hline
\multirow{1}{2.8cm}{\textbf{Luồng sự kiện thay thế}} &
\multicolumn{3}{c|}{%
\begin{tabular}{|p{0.8cm}|p{2.8cm}|p{6.2cm}|}
\hline
\textbf{STT} & \textbf{Thực hiện bởi} & \textbf{Hành động} \\
\hline
5a. & Hệ thống & Nếu UID không hợp lệ hoặc không tìm thấy thời khoá biểu công khai tương ứng, hiển thị thông báo thời khoá biểu công khai không tồn tại. \\
\hline
\end{tabular}} \\
\hline
\end{longtable}

\section{Phân tích hành vi}
\label{section:2.5}
\begingroup
\raggedbottom

Phần này phân tích hành vi của các chức năng trọng tâm thông qua biểu đồ hoạt
động và biểu đồ trình tự. Biểu đồ hoạt động làm rõ trình tự thực hiện, các nhánh
xử lý và điều kiện chuyển tiếp trong từng quy trình. Biểu đồ trình tự mô tả sự
tương tác theo thời gian giữa tác nhân, đối tượng biên, đối tượng điều khiển và
các thực thể nghiệp vụ tham gia vào từng kịch bản của ca sử dụng. Nội dung phân
tích tập trung vào các nghiệp vụ phân công giảng dạy, xếp và điều chỉnh thời
khoá biểu thủ công, xếp thời khoá biểu tự động và tra cứu thời khoá biểu. Bên
cạnh đó, biểu đồ trình tự đăng nhập được trình bày để làm rõ quá trình xác thực
trước khi Người xếp thời khoá biểu sử dụng các chức năng quản lý của hệ thống.

\begin{figure}[H]
    \centering
    \includegraphics[width=0.9\textwidth,height=0.72\textheight,keepaspectratio]{Hinhve/SequenceDiagram/sequence_dang_nhap.png}
    \caption{Biểu đồ tuần tự ca sử dụng Đăng nhập}
    \label{fig:sequence_dang_nhap}
\end{figure}

Hình~\ref{fig:sequence_dang_nhap} mô tả quá trình Người xếp thời khoá biểu nhập
thông tin đăng nhập và gửi yêu cầu xác thực. Hệ thống kiểm tra tính đầy đủ của
dữ liệu, tra cứu tài khoản và đối chiếu mật khẩu. Nếu thông tin hợp lệ, hệ thống
tạo thông tin xác thực và chuyển người dùng tới trang thời khoá biểu; ngược lại,
giao diện hiển thị thông báo lỗi tương ứng.

\begin{figure}[H]
    \centering
    \includegraphics[width=0.95\textwidth]{Hinhve/SequenceDiagram/sequence_phan_cong_chu_nhiem.png}
    \caption{Biểu đồ tuần tự ca sử dụng Phân công giáo viên chủ nhiệm}
    \label{fig:sequence_phan_cong_chu_nhiem}
\end{figure}

Hình~\ref{fig:sequence_phan_cong_chu_nhiem} mô tả trình tự tương tác giữa Người
xếp thời khoá biểu, giao diện phân công, đối tượng điều khiển phân công và các
thực thể lớp học, giáo viên trong quá trình phân công giáo viên chủ nhiệm. Biểu
đồ thể hiện cả trường hợp lựa chọn giáo viên chủ nhiệm mới và trường hợp loại bỏ
giáo viên chủ nhiệm hiện tại của lớp.

\begin{figure}[H]
    \centering
    \includegraphics[width=0.95\textwidth]{Hinhve/SequenceDiagram/sequence_phan_cong_bo_mon.png}
    \caption{Biểu đồ tuần tự ca sử dụng Phân công giáo viên bộ môn}
    \label{fig:sequence_phan_cong_bo_mon}
\end{figure}

Hình~\ref{fig:sequence_phan_cong_bo_mon} mô tả trình tự tải dữ liệu phân công,
hiển thị danh sách theo chế độ xem được lựa chọn và lưu giáo viên phụ trách môn
học của từng lớp. Khi môn học đã có giáo viên phụ trách, phân công hiện tại được
thay thế bằng giáo viên mới. Sau khi lưu, hệ thống cập nhật tổng số tiết đã phân
công và hiển thị cảnh báo nếu số tiết của giáo viên vượt định mức tham khảo.

\begin{figure}[H]
    \centering
    \includegraphics[width=0.95\textwidth,height=0.78\textheight,keepaspectratio]{Hinhve/SequenceDiagram/sequence_xep_tkb_thu_cong.png}
    \caption{Biểu đồ tuần tự ca sử dụng Xếp và điều chỉnh thời khoá biểu thủ công}
    \label{fig:sequence_xep_tkb_thu_cong}
\end{figure}

Hình~\ref{fig:sequence_xep_tkb_thu_cong} mô tả quá trình tải dữ liệu của tuần
được chọn, đối chiếu các tiết đã xếp với thông tin phân công hiện tại và làm nổi
bật các tiết bị xung đột. Trong quá trình xếp, người dùng lựa chọn từng ô trên
lưới và chọn môn học cần bố trí, thay thế hoặc loại bỏ. Hệ thống kiểm tra các
ràng buộc về giáo viên và phòng chức năng trước khi cập nhật giao diện ở trạng
thái chưa lưu. Khi hoàn tất, người dùng có thể lưu riêng tuần đang chỉnh sửa hoặc
áp dụng kết quả từ tuần đó trở đi.

\begin{figure}[H]
    \centering
    \includegraphics[width=0.95\textwidth,height=0.75\textheight,keepaspectratio]{Hinhve/SequenceDiagram/sequence_xep_tkb_tu_dong.png}
    \caption{Biểu đồ tuần tự ca sử dụng Xếp thời khoá biểu tự động}
    \label{fig:sequence_xep_tkb_tu_dong}
\end{figure}

Hình~\ref{fig:sequence_xep_tkb_tu_dong} mô tả quá trình kiểm tra dữ liệu phân
công, tải thông tin tuần, các tiết đã xếp và phòng chức năng trước khi chuyển dữ
liệu cùng các ràng buộc cho bộ xếp lịch tự động. Bộ xếp lịch tạo phương án cho
các tiết còn thiếu và giữ nguyên những tiết đã được bố trí trước đó. Kết quả được
hiển thị ở trạng thái chưa lưu để người dùng xem xét, sau đó hệ thống lưu theo
phạm vi tuần được lựa chọn và cập nhật lại lưới thời khoá biểu.

\begin{figure}[H]
    \centering
    \includegraphics[width=0.95\textwidth,height=0.75\textheight,keepaspectratio]{Hinhve/SequenceDiagram/sequence_tra_cuu_tkb.png}
    \caption{Biểu đồ tuần tự ca sử dụng Tra cứu thời khoá biểu}
    \label{fig:sequence_tra_cuu_tkb}
\end{figure}

Hình~\ref{fig:sequence_tra_cuu_tkb} mô tả quá trình Người dùng vãng lai mở URL
công khai chứa UID của thời khoá biểu. Sau khi tìm thấy thời khoá biểu tương
ứng, hệ thống tải thông tin năm học, học kì và danh sách tuần, tự động chọn tuần
đầu tiên rồi lấy các tiết học của tuần đó để hiển thị trên lưới. Khi người dùng
chọn tuần khác, hệ thống chỉ tải lại danh sách tiết học của tuần mới. Nếu UID
không hợp lệ hoặc thời khoá biểu công khai không tồn tại, hệ thống không hiển
thị dữ liệu thời khoá biểu và thông báo cho người dùng về trạng thái này.

\endgroup
\section{Yêu cầu phi chức năng}
\label{section:2.6}

\subsection{Hiệu năng hệ thống}

Các chức năng quản lý dữ liệu, phân công giảng dạy và tra cứu thời khoá biểu cần
phản hồi trong thời gian không vượt quá 2 giây ở điều kiện tải thông thường.
Việc kiểm tra xung đột khi người dùng xếp hoặc điều chỉnh một tiết học cần được
thực hiện trong thời gian dưới 1 giây để không làm gián đoạn thao tác. Hệ thống
được xây dựng cho quy mô một trường tiểu học với khoảng 50 giáo viên, 30 lớp
học và 10 môn học. Đối với chức năng xếp thời khoá biểu tự động, hiệu năng
được đánh giá dựa trên thời gian xử lý với các bộ dữ liệu có quy mô khác nhau
về số lớp học, giáo viên, môn học và phân công giảng dạy; kết quả đo lường sẽ
được trình bày trong phần kiểm thử hệ thống.

\subsection{Tính bảo mật}

Mật khẩu của người dùng phải được mã hoá bằng BCrypt trước khi lưu vào cơ sở
dữ liệu và không được lưu dưới dạng văn bản thuần. Hệ thống sử dụng JWT để
xác thực; token được lưu trong cookie có thuộc tính \texttt{httpOnly} nhằm hạn
chế nguy cơ token bị truy cập bởi mã JavaScript phía trình duyệt. Trong môi
trường triển khai thực tế, toàn bộ giao tiếp giữa trình duyệt, frontend và backend
phải được thực hiện qua HTTPS. CORS chỉ chấp nhận yêu cầu từ các nguồn đã
được khai báo, đồng thời các API thực hiện thao tác quản lý dữ liệu chỉ cho phép
người dùng đã xác thực truy cập.

\subsection{Khả năng tương thích}

Hệ thống hoạt động trên trình duyệt web và không yêu cầu người dùng cài đặt
thêm phần mềm. Giao diện cần hiển thị ổn định trên các phiên bản phổ biến gần
đây của Chrome, Microsoft Edge, Firefox và Safari. Hệ thống được ưu tiên thiết
kế cho máy tính để bàn và máy tính xách tay, đồng thời bố cục có khả năng điều
chỉnh theo các kích thước màn hình thông dụng.

\subsection{Tính khả dụng và độ tin cậy}

Hệ thống cần duy trì khả năng truy cập ổn định trong điều kiện vận hành thông
thường và hiển thị thông báo rõ ràng khi yêu cầu không hợp lệ hoặc không thể
kết nối tới máy chủ. Cơ sở dữ liệu PostgreSQL sử dụng các ràng buộc khoá ngoại
để bảo đảm mối liên hệ giữa các dữ liệu. Những thao tác cập nhật nhiều dữ liệu
liên quan, đặc biệt là phân công giảng dạy và xếp thời khoá biểu, cần được xử lý
trong giao dịch nhằm hạn chế trạng thái không nhất quán khi xảy ra lỗi. Trong
môi trường triển khai thực tế, dữ liệu cần được sao lưu định kỳ và có khả năng
phục hồi khi xảy ra sự cố.

\subsection{Khả năng bảo trì và mở rộng}

Frontend và backend được tổ chức thành hai thành phần độc lập, có thể xây dựng
và triển khai riêng với điều kiện duy trì tính tương thích của hợp đồng API giữa
hai phía. Mã nguồn được phân chia theo các nhóm chức năng và tầng xử lý nhằm
hỗ trợ kiểm tra, sửa lỗi và bổ sung nghiệp vụ. Frontend sử dụng TypeScript để
phát hiện sớm các lỗi kiểu dữ liệu, trong khi cấu trúc cơ sở dữ liệu được quản lý
bằng Flyway để theo dõi các thay đổi qua từng phiên bản. Các API được xây dựng
theo phong cách RESTful, tạo điều kiện để ứng dụng khác tái sử dụng những chức
năng hiện có. Dữ liệu được tổ chức theo năm học, cho phép bổ sung năm học mới
mà không làm thay đổi dữ liệu của các năm học trước.

\subsection{Khả năng sử dụng}

Giao diện thời khoá biểu sử dụng cấu trúc bảng theo thứ và tiết quen thuộc, các
thao tác chính được trình bày trực quan nhằm giúp người dùng dễ tiếp cận. Hệ
thống cung cấp thông báo kết quả sau mỗi thao tác và yêu cầu xác nhận đối với
những hành động có thể làm thay đổi hoặc mất dữ liệu. Các vị trí xảy ra xung
đột và trường hợp giáo viên vượt định mức được thể hiện trực tiếp trên giao diện,
giúp người dùng nhanh chóng nhận biết và xử lý.


%%%%%%%%%%%%%%%%%%%%%%%%%%%%%%%%%%%

\end{document}
